% ============================================================================================
% BAB I PENDAHULUAN — KERANGKA (bullet), bukan prosa jadi.
%
% Sumber: DECISIONS.md, bagian "Bab I — two questions about what the thesis claims to be"
% (baris 226) dan tabel "Contents by chapter" (baris 113).
%
% Dua entri menyentuh bab ini, keduanya bersyarat: D-21 dan D-28. Keduanya hanya
% mengubah Bab I bila kalimat ruang lingkup atau judul berubah — lihat TODO(open)
% di bawah. DECISIONS.md menegaskan: SATU percakapan pembimbing, bukan dua.
%
% Catatan: Latar Belakang dan Tujuan di sini TIDAK menyalin proposal II4091.
% Komentar penguji terhadap proposal adalah "isi Pendahuluan seharusnya ada di
% bab lain dan sebaliknya", jadi kerangka ini dibangun ulang. Tujuan ditulis
% sebagai sasaran terukur (angka, ambang, prosedur), bukan "lebih baik" atau
% "lebih efisien".
%
% Tidak ada \autocite / \cite di berkas ini. daftar-pustaka.bib di paper/Rafly TA/
% masih berkas dummy template, sehingga setiap kunci akan salah atau — lebih buruk —
% diam-diam teratasi ke laudon2020. Kebutuhan sitasi ditulis sebagai "% CITE:".
% ============================================================================================
\cleardoublepage
\chapter{PENDAHULUAN}
\label{chap:pendahuluan}

\section{Latar Belakang}

\begin{itemize}
    \item Petir adalah penyebab gangguan yang signifikan pada sistem transmisi
          dan distribusi tenaga listrik; kerapatan sambaran adalah masukan
          langsung bagi perancangan proteksi.
          % CITE: standar proteksi petir yang memakai GFD sebagai masukan (mis. IEC 62305 / SPLN)
          % CITE: statistik gangguan saluran transmisi akibat petir di Indonesia
    \item Kerapatan sambaran petir ke tanah (\textit{ground flash density}, GFD)
          secara konvensional dinyatakan dalam sambaran per km\textsuperscript{2}
          per tahun.
          % CITE: definisi baku GFD dan satuannya
    \item PLN Puslitbang mengoperasikan \textit{Lightning Detection System} (LDS)
          yang merekam sambaran di Jawa Barat; data inilah yang menjadi sumber
          domain tropis pada TA ini (D-22).
    \item Studi terdahulu memodelkan GFD pada resolusi \textbf{bulanan}.
          % CITE: studi predecessor — pemodelan GFD bulanan, wilayah Jawa Barat
    \item Celah yang diisi TA ini: dua sumber prediktor utama (ERA5 dan NASA
          POWER) tersedia \textbf{per jam}, sehingga agregasi ke bulanan membuang
          resolusi yang sebenarnya dimiliki data. Satu baris pada TA ini adalah
          satu sel 0,5$^\circ$ untuk satu jam (D-21).
          % TODO(open): Pertanyaan terbuka 7 — argumen positif untuk resolusi jam
          %   BELUM PERNAH DITULIS di mana pun. Dalam DECISIONS.md alasannya
          %   tercatat hanya sebagai RECOLLECTION: "pemetaan masukan-keluaran yang
          %   lebih akurat antara keadaan meteorologi dan petir yang menyertainya".
          %   Karena resolusi jam menyimpang dari mitigasi kelangkaan data yang
          %   sudah disetujui di proposal, argumen ini harus disusun penuh —
          %   di Bab III, bukan di sini. Bab I hanya menyebut konsekuensinya.
    \item Konsekuensi resolusi jam: target menjadi sangat didominasi nol —
          \textbf{94,30\% tropis} dan \textbf{97,00\% subtropis} (D-20). Ini
          dibawa apa adanya dan dilaporkan, tidak direkayasa hilang.
    \item \textit{Quantum machine learning} (QML) diajukan sebagai kandidat model
          karena klaim ekspresivitas sirkuit tervariasi pada jumlah parameter
          kecil; klaim efisiensi sampel diuji, bukan diasumsikan (D-14).
          % CITE: variational quantum circuit / QNN sebagai model regresi
          % CITE: klaim efisiensi sampel pada QML
          % CITE: barren plateau dan mitigasinya lewat fungsi biaya lokal
    \item Dua domain iklim dipakai agar generalisasi dapat diukur, bukan
          diklaim: tropis (Jawa Barat, PLN LDS) dan subtropis (Florida, NASA
          MERLIN), keduanya 2018--2024.
\end{itemize}

% TODO(open): Pertanyaan terbuka 8 — apakah keluaran per sel-jam masih sah
%   disebut "ground flash density"? GFD konvensional adalah flash/km2/tahun;
%   hitungan per sel-jam adalah OKURENSI petir, yang punya literatur dan
%   metriknya sendiri (POD, FAR, CSI, Brier). Bila build per jam yang
%   dilaporkan, kalimat ruang lingkup di subbab ini — dan mungkin judul —
%   harus menyatakannya.
%   Settled by: percakapan pembimbing. GABUNG dengan pertanyaan 9 di bawah.
% TODO(open): Pertanyaan terbuka 9 — apakah "Framework Qiskit" pada judul
%   bertahan? Sirkuit dibangun dan diverifikasi terhadap Qiskit (sepakat
%   sampai 1e-10 pada forward pass), tetapi SETIAP gradien pada hasil yang
%   dilaporkan dihitung PennyLane lightning.qubit dengan metode adjoint
%   (D-28). Posisi ini dapat dipertahankan, tetapi hanya bila dinyatakan.
%   Settled by: percakapan pembimbing yang SAMA dengan pertanyaan 8.

\section{Rumusan Masalah}

\begin{itemize}
    \item Bagaimana membangun tabel prediktor--target beresolusi sel-jam yang
          konsisten untuk dua domain iklim dari empat sumber yang berbeda
          resolusi dan konvensi waktunya (PLN/MERLIN, NASA POWER, ERA5)?
    \item Bagaimana kinerja \textit{quantum neural network} (QNN) dibandingkan
          tangga model klasik pada paritas parameter, untuk target yang
          94--97\% bernilai nol?
    \item Sejauh mana model yang dilatih pada satu domain iklim dapat
          diterapkan pada domain lain, dan berapa besar selisih kinerjanya?
    \item Pada resolusi pelaporan mana (1, 3, 6, atau 24 jam) kebutuhan
          nonfungsional NF-01 dapat dievaluasi secara bermakna, dan apakah ia
          terpenuhi?
\end{itemize}

\section{Tujuan}

% Setiap tujuan di bawah punya kriteria lulus/gagal yang dapat diperiksa.
% Hindari rumusan "lebih baik", "lebih efisien", "lebih optimal" — tidak ada
% ambang yang bisa diuji di baliknya.

\begin{itemize}
    \item \textbf{T1.} Membangun dua tabel terproses beresolusi sel-jam untuk
          2018--2024. \textit{Kriteria:} tropis 1.841.040 baris (30 sel) dan
          subtropis 3.314.304 baris (56 sel), 13 prediktor dasar di kedua
          domain (D-01), \textbf{tanpa satu pun nilai hilang}, sehingga tidak
          ada langkah imputasi di mana pun.
    \item \textbf{T2.} Mengukur QNN terhadap tangga enam \textit{rung}
          (\texttt{baseline\_trivial}, \texttt{ridge}, \texttt{nn\_matched},
          \texttt{qnn}, \texttt{nn\_large}, \texttt{nn\_full}) pada paritas
          parameter \textbf{47 lawan 52} (D-34). \textit{Kriteria:} setiap
          angka dilaporkan sebagai rerata $\pm$ simpangan baku atas
          \textbf{3 fold} (D-06) $\times$ \textbf{3 seed} (D-29), dan setiap
          tabel metrik memuat skor \textit{baseline} trivial dalam satuan yang
          sama.
    \item \textbf{T3.} Melaporkan NF-01 (NRMSE $< 0,1$ dan R\textsuperscript{2}
          $> 0,9$) terhadap \textbf{setiap} jendela pelaporan 1/3/6/24 jam
          (D-04). \textit{Kriteria:} kurva jendela lengkap, disertai pangsa nol
          per jendela. NF-01 dinyatakan \textbf{di muka tidak akan terpenuhi}
          pada resolusi jam (D-12) — ini pernyataan, bukan hasil yang ditunggu.
    \item \textbf{T4.} Mengukur selisih kinerja lintas domain dengan
          \textit{scaler} domain sumber, bukan domain target (D-36).
          \textit{Kriteria:} empat kombinasi latih--uji dilaporkan pada metrik
          yang sama.
    \item \textbf{T5.} Memenuhi NF-02 (reprodusibilitas) dengan artefak yang
          disebutkan namanya: satu \texttt{requirements.txt} tersemat eksak
          (D-41), satu \texttt{RANDOM\_SEED = 18222067} yang diimpor dan tidak
          disalin, serta \textit{config snapshot} yang menyertai setiap hasil
          (D-42).
          % TODO(open): status NF-02 — D-42 menyatakan NF-02 terpenuhi SECARA
          %   MEKANISME tetapi BELUM DIUJI ujung-ke-ujung (pertanyaan terbuka
          %   28: replay penuh belum pernah dijalankan). Tulis persis demikian.
\end{itemize}

\section{Batasan Masalah}

\begin{itemize}
    \item Wilayah dan periode: Jawa Barat dan Florida, 2018-01-01 sampai
          2024-12-31, grid 0,5$^\circ$ dengan kotak dibulatkan ke luar
          (D-15).
    \item Domain tropis memakai sambaran \textit{cloud-to-ground} saja,
          terverifikasi di sumber untuk ketujuh tahun (D-22).
          % TODO(open): Pertanyaan terbuka 10 — komposisi CG/IC ekspor MERLIN
          %   BELUM diketahui; tidak ada kolom diskriminasi. Selama belum
          %   selesai, SETIAP hasil subtropis bersifat provisional dan harus
          %   dinyatakan demikian. Rinciannya di Bab III.
    \item Eksekusi kuantum adalah \textbf{simulasi} \textit{statevector}
          (\texttt{SHOTS = None}, nilai harap eksak), bukan perangkat keras
          kuantum (D-33); gradien dihitung PennyLane \texttt{lightning.qubit}
          (D-28).
    \item Indeks K ERA5 (\texttt{KX}) tidak dipakai model mana pun, sehingga
          kedua domain menyajikan 13 prediktor dasar yang identik (D-01).
    \item Anggaran komputasi membatasi ukuran latih: \texttt{TRAIN\_ROWS}
          = 3.000, \texttt{MAX\_EPOCHS} = 30, 3 seed. Lima seed tidak muat
          dalam satu malam (D-29).
    \item Data mentah PLN tidak boleh diredistribusi dan tidak disertakan dalam
          repositori (D-43).
\end{itemize}

\section{Metodologi}

\begin{itemize}
    \item Studi literatur: GFD dan jaringan deteksi petir, pemodelan target
          \textit{zero-inflated}, dan QML tervariasi (Bab~\ref{chap:studi-literatur}).
    \item Akuisisi empat sumber data dan pembangunan tabel sel-jam untuk kedua
          domain (Bab~\ref{chap:analisis-masalah}).
    \item Analisis kualitas, eksklusi, dan bentuk dataset; perumusan kebutuhan
          fungsional dan nonfungsional (Bab~\ref{chap:analisis-masalah}).
    \item Perancangan target, prapemrosesan, protokol eksperimen, dan
          implementasi kedua arm model (Bab~\ref{chap:perancangan}).
    \item Implementasi paket \texttt{gfd\_data} dan \texttt{gfd\_model}
          (Bab~\ref{chap:implementasi}).
    \item Evaluasi: tangga model, kurva jendela, sweep OFAT, dan uji lintas
          domain (Bab~\ref{chap:evaluasi}).
\end{itemize}

\section{Sistematika Penulisan}

\begin{itemize}
    \item \textbf{Bab I Pendahuluan} --- latar belakang, rumusan masalah,
          tujuan terukur, batasan, dan metodologi.
    \item \textbf{Bab II Studi Literatur} --- petir dan GFD, jaringan deteksi,
          sumber data meteorologi, pemodelan target jarang, dan QML.
    \item \textbf{Bab III Analisis Masalah} --- kondisi saat ini, sumber dan
          akuisisi data, kualitas dan eksklusi, bentuk dataset, kebutuhan, dan
          pemilihan solusi.
    \item \textbf{Bab IV Perancangan} --- target dan prapemrosesan, protokol
          eksperimen, dan rancangan implementasi kedua model.
    \item \textbf{Bab V Implementasi} --- realisasi paket \texttt{gfd\_data}
          dan \texttt{gfd\_model}.
    \item \textbf{Bab VI Evaluasi} --- metode evaluasi, hasil, dan pembahasan.
    \item \textbf{Bab VII Penutup} --- kesimpulan dan saran.
\end{itemize}
